\documentclass[10pt,a4paper]{amsart}
\usepackage[margin=19mm]{geometry}
\usepackage{amsmath,amssymb,mathtools}
\usepackage[scr=rsfs,cal=euler]{mathalfa}
\usepackage{xfrac}
\usepackage[unicode,hidelinks,bookmarks=false]{hyperref}
\newcommand{\ie}{i.e.\ }
\newcommand{\cf}{cf.\ }
\newcommand{\resp}{respectively\ }
\newcommand{\Subsection}{\subsection}
\providecommand{\nobreakdash}{\nobreak\hbox{-}}
\setlength{\emergencystretch}{2em}
\multlinegap0pt
\usepackage{dsfont}
\usepackage{amssymb}
\usepackage{mathtools}
\usepackage{caption}
\usepackage[shortlabels]{enumitem}
\usepackage{tikz}
\usepackage[all]{xy}
\usepackage{xcolor}
\newtheorem{proposition}{Proposition}
\DeclareMathOperator{\Diag}{Diag}
\newcommand{\Q}{\mathbb{Q}}
\newcommand\Qp{{\mathbb{Q}_p}}
\newcommand{\Z}{\mathbb{Z}}
\newcommand\Zp{{\mathbb{Z}_p}}
\newcommand\bP{\mathbb{P}}
\newcommand{\cyc}{\textrm{cyc}}
\newcommand\gp{\mathfrak{p}}
\newcommand{\ob}[1]{\mkern 1.5mu\overline{\mkern-1.5mu#1\mkern-1.5mu}\mkern 1.5mu}
\DeclareMathOperator{\rk}{rk}
\newcommand\sC{\mathscr{C}}
\newcommand\sZ{\mathscr{Z}}
\title{On the rank of Leopoldt's and Gross's regulator maps}
\author{Alexandre Maksoud}
\date{Source: arXiv:2201.08203v2 (CC BY 4.0)}
\begin{document}
\maketitle

\noindent\emph{Source and licence.} On the rank of Leopoldt's and Gross's regulator maps. Version arXiv:2201.08203v2 (CC BY 4.0), distributed by arXiv under the Creative Commons Attribution 4.0 International licence, http://creativecommons.org/licenses/by/4.0/, as recorded in the arXiv licence metadata for this exact version and shown on the abstract page https://arxiv.org/abs/2201.08203v2. Full licence notice: \texttt{licenses/arxiv-2201.08203-CC-BY-4.0.txt}.\par\smallskip

\noindent\textit{Editorial scope.} Counted unit: the article's proposition prop:non\_cyc\_trivial\_char (source0.tex lines 736--741). The article's complete original proof of it is delivered with it (source0.tex lines 742--759, one \texttt{proof} environment of 18 lines). No mathematics is written, repaired or re-derived: the statement and the proof are the article's own lines, in the article's own notation, and no retained line is substituted or re-flowed.  Retained uncounted as the source-local context the proof needs: lines 705--719.  Nothing was omitted from any retained span, so the package carries no omission record.  No retained line contains a citation, so the wrapper carries no bibliography and no bibliography entry.  No retained line contains a figure environment, an image-inclusion command or drawing code, so no image is delivered and the package declares no asset dependency.  Editorial formatting only, in addition: an AMS \texttt{amsart} preamble that restates the article's own declarations and macros used by the retained lines, namely the environments \texttt{proposition} on separate counters, together with the standard packages those lines need.  The retained environments print in this document as Proposition 1 in order of appearance, whereas the article prints the same items with the numbering of its own full document.  The complete native source of the article as distributed in the arXiv e-print for this exact version is bundled unchanged as \texttt{sources/121043/source0.tex}.
	\begin{proposition}\label{prop:prop_clef_non_cyc}
		\begin{enumerate}
			\item The map sending a $\Zp$-extension $k_\infty/k$ to its coordinates $S$ defines a bijection between $\sZ(k)$ and $\{S \in \bP^{n-1}(\Qp) \ \colon \ M_{\mathds{1}}S=0\}$.
			\item Let $k_\infty \in \sZ(k)$ with coordinates $S\in \ker(M_\mathds{1})$ and consider the matrix of size $(n+r(\varphi))\times n$ given in block notation by:
			\begin{equation}\label{eq:def_N_phi_S}
				\begin{bmatrix}
				N_\varphi(S)=\Diag(S)L_\varphi - \Diag(L_\mathds{1}S)\\
				M_\varphi
			\end{bmatrix},
			\end{equation}
			where $\Diag(U)$ denotes the diagonal matrix associated with a column matrix $U$.
			Then $$\delta_{k_\infty/k}^\textbf{G}(\varphi) = \dim \ker N_\varphi(S) - \gamma,$$ 
			where $\gamma=1$ if $\varphi=\mathds{1}$ and $\gamma=0$ otherwise. 
		\end{enumerate}
	\end{proposition}

	\begin{proposition}\label{prop:non_cyc_trivial_char}
		Assume that $n-r\leq 2$. \begin{enumerate}
			\item  Let $\sC(k)=\{k_\infty\in\sZ(k) \ \colon\ \delta_{k_\infty/k}^\textbf{G}\neq 0\}$. If $\sC(k)\neq \sZ(k)$, then $\sC(k)$ is finite with $|\sC(k)|\leq n-1$. If, moreover, $k$ is imaginary quadratic, then $\sC(k)=\emptyset$.
			\item Let  $\sC_\varphi(k)=\{k_\infty\in\sZ(k) \ \colon\ \delta_{k_\infty/k}^\textbf{G}(\varphi)\neq 0\}$. If $\sC_\varphi(k)\neq \sZ(k)$, then $\sC_\varphi(k)$ is finite with $|\sC(k)|\leq n$.
		\end{enumerate}
	\end{proposition}

	\begin{proof}
		Let us prove (1). Notice first that all $S\in \ker M$ lie in the kernel of $N(S)$. By Proposition \ref{prop:prop_clef_non_cyc} with $\varphi=\mathds{1}$ and $K=k$, $\sC(k)$ is in bijection with the set of all $S\in \bP^{n-1}(\Qp)$ such that $MS=0$ and $\rk N(S)<n-1$.
		
		Assume first that $k$ is quadratic, so $n=2$ and $r=0$. Given any $S=(s_1,s_2)\in \bP^1(\Qp)$, the matrix $N(S)=\Diag(S)L - \Diag(LS)$ has the form 
		\[\begin{pmatrix}
			-s_2L_{1,2} & s_1 L_{1,2} \\ s_2 L_{2,1} & -s_1 L_{2,1}
		\end{pmatrix}.\]
		As $\log_p$ is injective on $\Z_p^\times$, $\log_p(\iota_{\gp_1}(u_2))$ and $\log_p(\iota_{\gp_2}(u_1))$ both are non-zero, so at least one of the two non-diagonal entries of $N(S)$ is non-zero. Therefore, this matrix has rank one for any $S\in \bP^1(\Qp)$, and $\sC(k)=\emptyset$.
		
		We no longer assume that $k$ is imaginary quadratic, but we still assume that $n-r\leq 2$. The case where $n-r=1$ is trivial, because it forces $\sZ(k)=\{k_{\cyc}\}$. We may then assume that $n-r=2$. Since $M$ has full rank $r$, there exist invertible matrices $P,Q$ such that $PMQ=\left( I_r\ | \ 0\right)$, where $I_r$ is the identity matrix of size $r$. The change of variables $S'=Q^{-1}S$ induces a linear bijection between $\ker M \subset \bP^{n-1}(\ob{\Q}_p)$ and the projective line $\{0\} \times \bP^{1}(\ob{\Q}_p) \subset \bP^{n-1}(\ob{\Q}_p)$. Now consider the list $P_1(S'),\ldots,P_t(S')$ of all $(n-1)\times(n-1)$-minors of the matrix	
		\[N(S)=N(QS')=\begin{bmatrix}
			\Diag(QS')L- \Diag(LQS')\\
			M
		\end{bmatrix}.\]
		All the $P_k$'s are two-variable homogeneous polynomials of degree $\leq n-1$. In particular, if $\sC(k)\neq \sZ(k)$, then at least one of the $P_k$'s is not the zero polynomial and hence, it has at most $n-1$ zeros in $\{0\} \times \bP^{1}(\ob{\Q}_p)$, so we can conclude that $|\sC(k)|\leq n-1$.
		
		The proof of point (2) very similar to the previous one. Indeed, by Proposition \ref{prop:prop_clef_non_cyc}, $\sC_\varphi(k)$ is in bijection with the set of all $S\in \bP^{n-1}(\Qp)$ such that $MS=0$ and $\rk N_\varphi(S)<n$. The same argument with the $(n-1)\times(n-1)$ minors of $N(S)$ replaced by the $n\times n$ minors of $N_\varphi(S)$ shows that, if $\sC_\varphi(k)\neq \sZ(k)$, then $|\sC(k)|\leq n$.
	\end{proof}

\end{document}
